% ============================================================
%  Standalone Literature Review
%  Topic: GNN-Based Fault Detection and Localization in Power Grids
% ============================================================
\documentclass[12pt,a4paper]{article}

\usepackage[T1]{fontenc}
\usepackage[utf8]{inputenc}
\usepackage{amsmath,amssymb}
\usepackage{booktabs}
\usepackage{array}
\usepackage{geometry}
\usepackage{setspace}
\usepackage{cite}
\usepackage{hyperref}
\usepackage{microtype}

\geometry{margin=1in}
\onehalfspacing

\hypersetup{
    colorlinks=true,
    linkcolor=blue,
    citecolor=blue,
    urlcolor=blue,
    pdftitle={Literature Review on GNN-Based Fault Detection and Localization in Power Grids},
    pdfauthor={Anik Kumar Basu, Sabrina Shonda Snaha, Md. Sadik Mahmud Adive}
}

\title{\textbf{Literature Review on Graph Neural Networks for Fault Detection and Localization in Power Grids}}
\author{
Anik Kumar Basu \and
Sabrina Shonda Snaha \and
Md. Sadik Mahmud Adive\\
Department of Computer Science and Engineering\\
Bangladesh University of Business and Technology, Bangladesh
}
\date{June 2026}

\begin{document}

\maketitle

\begin{abstract}
Fault detection and localization are essential for secure and reliable smart-grid operation. As power systems integrate renewable energy sources, bidirectional power flows, distributed generation, and high-frequency sensing devices such as phasor measurement units (PMUs), conventional protection and monitoring methods face increasing difficulty. This literature review examines major research directions for power-system fault detection and localization, including traditional relay-based methods, shallow machine learning, deep learning, graph neural networks (GNNs), and spatiotemporal learning approaches. The review shows that topology-agnostic methods often fail to represent the physical structure of power networks, while static graph methods may underuse the temporal evolution of fault transients. Spatiotemporal GNNs provide a promising direction because they combine graph-based spatial learning with temporal sequence modeling. The review concludes by identifying key research gaps related to bus-level localization, robustness, real-data validation, and multi-task learning.
\end{abstract}

\noindent\textbf{Keywords:} literature review, graph neural networks, smart grid, fault detection, fault localization, PMU, GCN, LSTM, spatiotemporal learning.

% ============================================================
\section{Introduction}
% ============================================================

Modern power grids are becoming more complex because of renewable energy integration, changing load behavior, distributed energy resources, and high-speed measurement systems. These changes make fault detection and localization more challenging. A protection system must not only detect whether a fault has occurred, but also identify where the fault originated so that corrective action can be taken quickly.

The literature on this problem has evolved through several stages. Early approaches used impedance-based relays, threshold logic, and signal-processing techniques. Later, researchers introduced shallow machine-learning models such as support vector machines (SVMs), random forests, and artificial neural networks (ANNs). More recent studies use deep learning, including convolutional neural networks (CNNs), recurrent neural networks, and GNNs.

The central challenge is that power-system faults are both spatial and temporal. They are spatial because disturbances propagate through buses, lines, transformers, and electrical coupling. They are temporal because voltage, current, frequency, and power signals evolve rapidly during pre-fault, fault, and post-fault periods. A strong fault-localization model should therefore represent both the grid topology and the transient signal behavior.

This review presents the main research directions and explains why spatiotemporal GNNs are a strong candidate for real-time fault detection and localization in power grids.

% ============================================================
\section{Traditional Fault Detection and Localization}
% ============================================================

Traditional fault detection methods include impedance-based relaying, overcurrent protection, differential protection, traveling-wave analysis, and threshold-based monitoring. These methods are widely used because they are fast, interpretable, and suitable for real-time protection hardware. Their decisions are also easier for system operators to understand because they are based on well-known electrical principles.

However, traditional methods depend heavily on accurate system parameters, relay settings, line impedance values, and stable operating assumptions. In modern grids, these assumptions are less reliable. Renewable generation, inverter-based resources, topology changes, and bidirectional power flow can alter fault signatures and reduce the reliability of fixed rules.

Liao et al. \cite{liao2021} discuss how smart-grid operation requires methods that can process large volumes of measurement data and adapt to changing system behavior. Chen et al. \cite{chen2020} also emphasize that data-driven fault diagnosis can address limitations of purely rule-based systems. Traditional methods remain important, but they are often insufficient for flexible, high-dimensional, PMU-based fault localization.

% ============================================================
\section{Shallow Machine-Learning Approaches}
% ============================================================

Shallow machine-learning methods were introduced to improve adaptability. Common examples include SVMs, decision trees, random forests, k-nearest neighbors, and shallow ANNs. These methods can learn nonlinear decision boundaries from measured voltage, current, power, or frequency features.

The main advantage of shallow learning is that it can outperform simple threshold rules when the training data are representative. For example, a classifier can learn patterns associated with different fault types or approximate fault locations. These models are also relatively easy to train compared with large neural networks.

The major limitation is feature dependence. Shallow models usually require manually designed features such as RMS values, harmonic components, sequence components, wavelet features, or statistical summaries. If the operating condition changes, those handcrafted features may no longer separate fault classes reliably. In addition, shallow models usually treat bus measurements as ordinary feature vectors and do not explicitly represent the physical grid topology.

For bus-level fault localization, ignoring topology is a serious limitation. Two buses may be electrically close even if their feature columns are far apart in a vector representation. This motivates graph-based learning methods that directly encode the power-network structure.

% ============================================================
\section{Deep Learning for Fault Diagnosis}
% ============================================================

Deep learning reduces the need for manual feature extraction. CNNs can learn patterns from waveform images, time-frequency maps, or bus-feature matrices. Recurrent models such as LSTMs can learn temporal patterns from sequential PMU windows.

CNN-based methods are useful because they can automatically detect local patterns in data. However, Owerko et al. \cite{owerko2020} point out that power systems are not naturally arranged as Euclidean grids. In fault localization, this matters because CNNs depend on the ordering of input rows and columns. If bus ordering changes, the physical system remains the same, but the CNN input pattern changes.

LSTM networks, introduced by Hochreiter and Schmidhuber \cite{hochreiter1997}, are useful for sequential data. In power-system fault analysis, LSTMs can capture how signals evolve over short time windows. This is important because early transient behavior often contains useful information about fault type and location.

However, a pure LSTM over flattened bus measurements does not know which buses are connected. It can learn time dynamics, but it does not explicitly model grid topology. Therefore, deep learning methods still need a way to combine temporal learning with physical network structure.

% ============================================================
\section{Graph Neural Networks in Power Systems}
% ============================================================

GNNs are well suited for power systems because a grid can naturally be represented as a graph. Buses are nodes, transmission lines and transformers are edges, and measurements such as voltage magnitude, phase angle, active power, reactive power, and frequency deviation are node features.

Kipf and Welling \cite{kipf2017} introduced a graph convolutional framework that updates node representations by aggregating information from neighboring nodes. In a power grid, this is meaningful because disturbances propagate through electrical connections. A GNN can therefore learn from physically connected buses rather than from arbitrary feature ordering.

Chen et al. \cite{chen2020} applied graph convolutional networks to power-system fault diagnosis and showed that topology-aware learning can improve diagnosis performance. Donon et al. \cite{donon2019} explored graph neural solvers for power-system problems and showed that graph learning can capture structural properties of electrical networks. James et al. \cite{james2020} investigated GCNs for fault location in distribution networks, showing the relevance of graph learning for localization tasks.

The strength of GNNs is physical consistency. If the graph representation is constructed correctly, the model can learn from electrical neighborhoods instead of relying on arbitrary bus indexing. GNNs also scale well for sparse networks, since real power grids have far fewer lines than all possible bus-to-bus connections.

The limitation is that many GNN-based approaches use static snapshots. A single post-fault snapshot may not capture enough information to distinguish the true fault location, especially when neighboring buses experience similar disturbances. This leads naturally to spatiotemporal graph learning.

% ============================================================
\section{Spatiotemporal Learning}
% ============================================================

Spatiotemporal learning combines spatial modeling with temporal modeling. Yu et al. \cite{yu2018} proposed a spatiotemporal graph convolutional network for traffic forecasting, showing that graph structure and temporal dynamics can be learned jointly. Although traffic networks and power grids are different, both involve signals that evolve over time on a network.

For power-system fault localization, the same idea is highly relevant. A GCN can model how the fault disturbance spreads across connected buses, while an LSTM or temporal convolution can model how the disturbance evolves across a PMU window. This combination is especially useful for real-time localization because the early transient may contain the strongest location-specific information.

A hybrid GCN-LSTM structure is therefore a practical architecture for this problem. The GCN component handles spatial propagation, and the LSTM component handles temporal progression. A multi-task objective can further improve learning by training the model to predict both fault location and fault type.

% ============================================================
\section{Comparison of Method Families}
% ============================================================

Table \ref{tab:comparison} summarizes the major method families reviewed in this document.

\begin{table}[h]
\centering
\small
\caption{Comparison of Method Families for Power-Grid Fault Analysis}
\label{tab:comparison}
\begin{tabular}{@{}p{0.22\textwidth}p{0.30\textwidth}p{0.34\textwidth}@{}}
\toprule
\textbf{Method Family} & \textbf{Strength} & \textbf{Limitation} \\
\midrule
Relay and threshold methods & Fast, interpretable, practical for protection systems & Sensitive to parameter uncertainty and changing grid conditions \\
Shallow machine learning & Learns nonlinear patterns from measured data & Requires handcrafted features and often ignores topology \\
CNN-based deep learning & Learns features automatically from waveform or matrix inputs & Assumes Euclidean structure and depends on bus ordering \\
LSTM/recurrent models & Captures temporal evolution of fault signals & Does not explicitly model physical grid connectivity \\
GNN/GCN models & Uses grid topology directly through message passing & Often applied to static snapshots without temporal context \\
Spatiotemporal GNNs & Models both grid topology and time-varying transients & Requires careful dataset design and robustness validation \\
\bottomrule
\end{tabular}
\end{table}

% ============================================================
\section{Research Gap}
% ============================================================

The reviewed literature reveals several research gaps. First, many studies emphasize fault detection or fault-type classification, while bus-level localization remains more difficult. Second, topology-agnostic models may perform well on fixed datasets but lack physical consistency under bus reordering or topology changes. Third, spatial GNN models often underuse temporal fault signatures. Fourth, many studies report accuracy but provide limited robustness analysis under measurement noise, class imbalance, and N-1 contingencies.

These gaps suggest the need for a model that jointly learns topology and time. A spatiotemporal GNN can address this by using graph convolution for spatial dependencies and sequence modeling for transient behavior. Multi-task learning can also help because fault type and fault location are related outputs.

% ============================================================
\section{Direction for Current Research}
% ============================================================

The reviewed works support a research direction based on three ideas. First, fault signals are graph-structured because physical connections determine how disturbances propagate. Second, fault signals are time-dependent because useful localization information appears during short transient windows. Third, fault location and fault type are related, so a shared representation can support both tasks.

For this reason, a model that combines GCN layers, LSTM layers, and multi-task output heads is well aligned with the limitations found in existing work. Synthetic data generation tools such as Pandapower \cite{thurner2018} can support systematic experiments by varying bus location, fault type, fault resistance, and loading level.

% ============================================================
\section{Conclusion}
% ============================================================

This literature review shows a clear progression from traditional protection methods to topology-aware spatiotemporal learning. Traditional and shallow-learning methods remain useful, but they struggle with high-dimensional PMU data, changing grid conditions, and bus-level localization. CNN and LSTM methods improve feature learning but do not fully solve the topology-time coupling problem. GNNs provide a natural way to represent power grids, while spatiotemporal GNNs extend this advantage by modeling transient dynamics.

Overall, the literature supports the use of spatiotemporal graph neural networks for real-time fault detection and localization. Future research should focus on real-data validation, robustness under noise and topology changes, explainability, and transfer learning across different grid systems.

% ============================================================
\begin{thebibliography}{9}

\bibitem{liao2021}
W.~Liao, B.~Bak-Jensen, J.~R.~Pillai \emph{et al.},
``Review of Graph Neural Networks for Smart Grids: An Overview,''
\emph{IEEE Power \& Energy Society General Meeting (PESGM)}, 2021.

\bibitem{owerko2020}
D.~Owerko, F.~Gama, and A.~Ribeiro,
``Optimal Power Flow using Graph Neural Networks,''
\emph{IEEE Int.\ Conf.\ Acoustics, Speech and Signal Processing (ICASSP)}, 2020.

\bibitem{chen2020}
K.~Chen, J.~Hu, and Y.~Zhang,
``Power System Fault Diagnosis Based on Graph Convolutional Networks,''
\emph{IEEE Int.\ Conf.\ Energy Internet and Energy System Integration (EI2)}, 2020.

\bibitem{donon2019}
B.~Donon, B.~Donnot, I.~Guyon, and M.~Sebag,
``Graph Neural Solver for Power Systems,''
\emph{Int.\ Joint Conf.\ Neural Networks (IJCNN)}, 2019.

\bibitem{yu2018}
B.~Yu, H.~Yin, and Z.~Zhu,
``Spatiotemporal Graph Convolutional Networks: A Deep Learning Framework
for Traffic Forecasting,''
\emph{IJCAI}, 2018.

\bibitem{james2020}
J.~James, Y.~C.~Chen, and A.~Dominguez-Garcia,
``Fault Location in Power Distribution Networks via Graph Convolutional Networks,''
\emph{IEEE North American Power Symposium (NAPS)}, 2020.

\bibitem{kipf2017}
T.~N.~Kipf and M.~Welling,
``Semi-Supervised Classification with Graph Convolutional Networks,''
\emph{Int.\ Conf.\ Learning Representations (ICLR)}, 2017.

\bibitem{hochreiter1997}
S.~Hochreiter and J.~Schmidhuber,
``Long Short-Term Memory,''
\emph{Neural Computation}, vol.~9, no.~8, pp.~1735--1780, 1997.

\bibitem{thurner2018}
L.~Thurner \emph{et al.},
``Pandapower---An Open-Source Python Tool for Convenient Modeling, Analysis,
and Optimization of Electric Power Systems,''
\emph{IEEE Trans.\ Power Systems}, vol.~33, no.~6, pp.~6510--6521, 2018.

\end{thebibliography}

\end{document}
